\section{Potencial eléctrico}

En el capítulo anterior estudiamos la energía potencial eléctrica \(U\), que depende
de la configuración de las cargas que forman un sistema.

Ahora introduciremos una magnitud relacionada, pero conceptualmente diferente:
el \textbf{potencial eléctrico}.

El potencial eléctrico permite describir energéticamente un campo eléctrico sin
necesidad de especificar qué carga se colocará en él.

Así como el campo eléctrico representa fuerza por unidad de carga,

\[
E=\frac{F}{q},
\]

el potencial eléctrico representa energía potencial por unidad de carga.


\subsection{Definición de potencial eléctrico}

Consideremos una carga de prueba \(q\) situada en un punto donde su energía potencial
eléctrica es \(U\).

El potencial eléctrico en ese punto se define como

\begin{equation}
\boxed{
V=\frac{U}{q}
}.
\end{equation}

Por tanto,

\begin{equation}
\boxed{
U=qV
}.
\end{equation}

\begin{definicion}[Potencial eléctrico]

El potencial eléctrico \(V\) en un punto es la energía potencial eléctrica por unidad
de carga que tendría una carga colocada en ese punto.

\end{definicion}

El potencial eléctrico es una \textbf{magnitud escalar}.

Por tanto, posee valor y signo, pero no dirección.


\begin{idea}

El campo eléctrico y el potencial describen aspectos diferentes de una misma
interacción:

\[
\boxed{\text{campo eléctrico}\longrightarrow\text{efecto de fuerza por unidad de carga}}
\]

\[
\boxed{\text{potencial eléctrico}\longrightarrow\text{energía por unidad de carga}}.
\]

\end{idea}


\subsection{El potencial es una propiedad del campo}

El valor de \(V\) en un punto está determinado por las cargas que producen el campo
y por la posición del punto.

No depende de la carga de prueba que posteriormente coloquemos allí.

Si en un mismo punto colocamos distintas cargas \(q\), cada una puede adquirir una
energía potencial diferente,

\[
U=qV,
\]

pero el potencial del punto permanece igual.

Por ello debe distinguirse cuidadosamente entre

\[
\boxed{V}
\]

y

\[
\boxed{U}.
\]

El potencial caracteriza el espacio creado por las cargas fuente, mientras que la
energía potencial depende además de la carga que interactúa con ese campo.


\section{Unidad del potencial eléctrico}

A partir de

\[
V=\frac{U}{q},
\]

la unidad SI del potencial es

\[
[V]=\frac{\si{\joule}}{\si{\coulomb}}.
\]

Esta unidad recibe el nombre de \textbf{voltio}:

\begin{equation}
\boxed{
\SI{1}{\volt}
=
\SI{1}{\joule\per\coulomb}
}.
\end{equation}

Esto significa que si una carga de \(\SI{1}{\coulomb}\) situada en determinado punto
posee una energía potencial de \(\SI{1}{\joule}\), el potencial en ese punto es
\(\SI{1}{\volt}\).


\section{Diferencia de potencial}

En muchas situaciones físicas no interesa el valor absoluto del potencial, sino su
cambio entre dos puntos.

Si una carga pasa de un punto \(A\) a un punto \(B\), definimos la diferencia de
potencial como

\begin{equation}
\boxed{
\Delta V=V_B-V_A
}.
\end{equation}

Como

\[
U=qV,
\]

se obtiene

\begin{equation}
\boxed{
\Delta U=q\Delta V
}.
\end{equation}

Por tanto,

\begin{equation}
\boxed{
\Delta V=\frac{\Delta U}{q}
}.
\end{equation}


\subsection{Diferencia de potencial y trabajo eléctrico}

Sabemos que

\[
\Delta U=-W_e.
\]

Por ello,

\begin{equation}
\boxed{
\Delta V=-\frac{W_e}{q}
}.
\end{equation}

Equivalentemente,

\begin{equation}
\boxed{
W_e=q(V_A-V_B)
}.
\end{equation}

Esta expresión indica que el trabajo realizado por la fuerza eléctrica depende de la
diferencia de potencial entre las posiciones inicial y final.


\subsection{Trabajo externo}

Si una carga se desplaza lentamente mediante una fuerza externa, sin cambio apreciable
de energía cinética,

\[
W_{\text{ext}}=\Delta U.
\]

Entonces,

\begin{equation}
\boxed{
W_{\text{ext}}=q\Delta V
}.
\end{equation}


\section{Significado físico de la diferencia de potencial}

Supongamos que una carga positiva se mueve espontáneamente bajo la acción exclusiva
de la fuerza eléctrica.

Como el campo realiza trabajo positivo,

\[
W_e>0,
\]

entonces

\[
\Delta U<0.
\]

Para una carga positiva,

\[
\Delta U=q\Delta V,
\]

por lo que

\[
\Delta V<0.
\]

Así, una carga positiva tiende espontáneamente a desplazarse hacia regiones de menor
potencial.

\begin{equation}
\boxed{
q>0
\quad\Longrightarrow\quad
\text{movimiento espontáneo hacia menor }V
}
\end{equation}

Para una carga negativa ocurre lo contrario.

Como

\[
U=qV,
\]

una carga negativa puede disminuir su energía potencial desplazándose hacia regiones
de mayor potencial.

\begin{equation}
\boxed{
q<0
\quad\Longrightarrow\quad
\text{movimiento espontáneo hacia mayor }V
}
\end{equation}

\begin{advertencia}

Las cargas no se mueven necesariamente hacia menor potencial.

La dirección del movimiento espontáneo depende del signo de la carga.

Lo que siempre disminuye en un movimiento espontáneo puramente electrostático es la
energía potencial \(U\).

\end{advertencia}


\section{Voltaje}

La expresión \textbf{voltaje} se utiliza frecuentemente como sinónimo de
\textbf{diferencia de potencial}.

Por tanto, cuando se afirma que entre dos puntos existe un voltaje de
\(\SI{12}{\volt}\), se está indicando una diferencia de potencial de

\[
\Delta V=\SI{12}{\volt}.
\]

El término potencial eléctrico suele referirse al valor \(V\) asociado con un punto,
mientras que voltaje suele emplearse para una diferencia \(V_B-V_A\).


\section{Potencial y campo eléctrico}

La diferencia de potencial puede calcularse directamente a partir del campo eléctrico.

Como

\[
W_e
=
q\int_A^B
\left(
E_x\,dx+E_y\,dy+E_z\,dz
\right)
\]

y

\[
\Delta V=-\frac{W_e}{q},
\]

obtenemos

\begin{equation}
\boxed{
V_B-V_A
=
-\int_A^B
\left(
E_x\,dx+E_y\,dy+E_z\,dz
\right)
}.
\end{equation}

Ésta es una de las relaciones fundamentales entre campo eléctrico y potencial.


\subsection{Interpretación del signo negativo}

El signo negativo indica que el potencial disminuye en el sentido del campo eléctrico.

Por tanto,

\begin{equation}
\boxed{
\text{el campo eléctrico apunta hacia potenciales decrecientes}
}.
\end{equation}

Si nos desplazamos en el sentido del campo,

\[
\Delta V<0.
\]

Si nos desplazamos en sentido contrario al campo,

\[
\Delta V>0.
\]


\section{Diferencia de potencial en un campo uniforme}

Consideremos un campo eléctrico uniforme.

Como sus componentes son constantes,

\[
V_B-V_A
=
-
\left(
E_x\Delta x+
E_y\Delta y+
E_z\Delta z
\right).
\]

Por tanto,

\begin{equation}
\boxed{
\Delta V
=
-
\left(
E_x\Delta x+
E_y\Delta y+
E_z\Delta z
\right)
}.
\end{equation}

Si el campo sólo tiene componente en \(x\),

\begin{equation}
\boxed{
\Delta V=-E_x\Delta x
}.
\end{equation}


\subsection{Desplazamiento en el sentido del campo}

Si

\[
E_x>0
\]

y

\[
\Delta x>0,
\]

entonces

\[
\Delta V=-E_x\Delta x<0.
\]

El potencial disminuye.


\subsection{Desplazamiento contrario al campo}

Si

\[
E_x>0
\]

pero

\[
\Delta x<0,
\]

entonces

\[
\Delta V>0.
\]

El potencial aumenta.


\subsection{Desplazamiento perpendicular al campo}

Si el desplazamiento es perpendicular al campo,

\[
E_x\Delta x+E_y\Delta y+E_z\Delta z=0.
\]

Entonces,

\begin{equation}
\boxed{\Delta V=0}.
\end{equation}

Los puntos inicial y final se encuentran al mismo potencial.


\subsection{Ejemplo: diferencia de potencial en un campo uniforme}

\begin{ejemplo}

En una región existe un campo eléctrico uniforme cuyas componentes son

\[
E_x=\SI{300}{\newton\per\coulomb},
\qquad
E_y=-\SI{200}{\newton\per\coulomb}.
\]

Determine la diferencia de potencial entre

\[
A=(0,0)
\]

y

\[
B=(\SI{0.40}{\metre},\SI{0.30}{\metre}).
\]

\end{ejemplo}

Los desplazamientos son

\[
\Delta x=\SI{0.40}{\metre},
\qquad
\Delta y=\SI{0.30}{\metre}.
\]

Aplicamos

\[
\Delta V
=
-(E_x\Delta x+E_y\Delta y).
\]

Sustituyendo,

\[
\Delta V
=
-\left[
(300)(0.40)+(-200)(0.30)
\right].
\]

Entonces,

\[
\Delta V=-(120-60).
\]

Por tanto,

\[
\boxed{
V_B-V_A=-\SI{60}{\volt}
}.
\]

El punto \(B\) se encuentra a un potencial \(\SI{60}{\volt}\) menor que el punto
\(A\).


\section{Potencial eléctrico de una carga puntual}

Consideremos una carga puntual fuente \(Q\).

El campo eléctrico a una distancia \(r\) tiene magnitud

\[
E=k\frac{|Q|}{r^2}.
\]

Para una distribución localizada de carga se adopta normalmente como referencia

\[
V(\infty)=0.
\]

Al integrar el campo de la carga puntual desde el infinito hasta una distancia \(r\),
se obtiene

\begin{equation}
\boxed{
V=k\frac{Q}{r}
}.
\end{equation}

En esta expresión el signo de \(Q\) debe conservarse.


\subsection{Carga positiva}

Si

\[
Q>0,
\]

entonces

\[
\boxed{V>0}.
\]

El potencial disminuye desde valores positivos cercanos a la carga hasta aproximarse
a cero cuando \(r\to\infty\).


\subsection{Carga negativa}

Si

\[
Q<0,
\]

entonces

\[
\boxed{V<0}.
\]

El potencial es muy negativo cerca de la carga y aumenta hacia cero al alejarnos.


\subsection{Dependencia con la distancia}

En ambos casos,

\[
\boxed{
|V|\propto\frac{1}{r}
}.
\]

Por tanto,

\[
r\to2r
\quad\Longrightarrow\quad
V\to\frac{V}{2}.
\]

Esto debe distinguirse de la dependencia del campo eléctrico:

\[
E\propto\frac{1}{r^2}.
\]


\subsection{Ejemplo: potencial de una carga puntual}

\begin{ejemplo}

Una carga puntual

\[
Q=+\SI{3.0}{\micro\coulomb}
\]

se encuentra en el origen.

Determine el potencial eléctrico a una distancia

\[
r=\SI{0.40}{\metre},
\]

tomando

\[
V(\infty)=0.
\]

\end{ejemplo}

Aplicamos

\[
V=k\frac{Q}{r}.
\]

Entonces,

\[
V
=
(8.99\times10^9)
\frac{3.0\times10^{-6}}{0.40}.
\]

Por tanto,

\[
\boxed{
V\approx\SI{6.7e4}{\volt}
}.
\]

El potencial es positivo porque la carga fuente es positiva.


\section{Potencial debido a varias cargas puntuales}

El potencial eléctrico obedece el principio de superposición.

Si varias cargas puntuales producen potencial en un punto \(P\),

\begin{equation}
\boxed{
V(P)=\sum_iV_i(P)
}.
\end{equation}

Como

\[
V_i=k\frac{q_i}{r_i},
\]

se obtiene

\begin{equation}
\boxed{
V(P)=k\sum_i\frac{q_i}{r_i}
}.
\end{equation}

donde \(r_i\) es la distancia desde la carga \(q_i\) hasta el punto \(P\).

\begin{idea}

La superposición del potencial es una suma \textbf{algebraica de escalares}.

No es necesario descomponer el potencial en componentes rectangulares.

\end{idea}

Los signos de las cargas sí deben conservarse:

\[
q_i>0\Longrightarrow V_i>0,
\]

\[
q_i<0\Longrightarrow V_i<0.
\]


\section{Campo y potencial: dos tipos distintos de superposición}

Para varias cargas puntuales, el campo requiere sumar sus componentes:

\[
E_x=\sum_iE_{i,x},
\qquad
E_y=\sum_iE_{i,y},
\]

mientras que el potencial se obtiene directamente mediante

\[
V=\sum_iV_i.
\]

Por ello, calcular el potencial suele ser algebraicamente más sencillo que calcular
el campo eléctrico.

\begin{advertencia}

Que el potencial sea escalar no significa que deban ignorarse los signos de las
cargas.

Se suman escalares, pero cada término

\[
k\frac{q_i}{r_i}
\]

puede ser positivo o negativo.

\end{advertencia}


\subsection{Ejemplo: potencial producido por dos cargas}

\begin{ejemplo}

Una carga

\[
q_1=+\SI{3.0}{\micro\coulomb}
\]

se encuentra en el origen y una carga

\[
q_2=-\SI{4.0}{\micro\coulomb}
\]

se encuentra en

\[
(\SI{0.40}{\metre},0).
\]

Determine el potencial eléctrico en

\[
P=(0,\SI{0.30}{\metre}),
\]

tomando \(V(\infty)=0\).

\end{ejemplo}

La situación es

\begin{center}
\begin{tikzpicture}[>=Latex,scale=1.05]

\draw[->] (-0.8,0) -- (5.2,0) node[right] {$x$};
\draw[->] (0,-0.7) -- (0,4.4) node[above] {$y$};

\filldraw (0,0) circle (2.2pt)
node[below left=3pt] {$q_1$};

\filldraw (4,0) circle (2.2pt)
node[below=4pt] {$q_2$};

\filldraw (0,3) circle (2.2pt)
node[left=4pt] {$P$};

\draw[dashed] (0,3) -- (4,0);

\node[left] at (0,1.5) {\(\SI{0.30}{\metre}\)};
\node[below] at (2,0) {\(\SI{0.40}{\metre}\)};
\node[above right] at (2,1.5) {\(\SI{0.50}{\metre}\)};

\end{tikzpicture}
\end{center}

La distancia desde \(q_1\) hasta \(P\) es

\[
r_1=\SI{0.30}{\metre}.
\]

Para \(q_2\),

\[
r_2=\sqrt{(0.40)^2+(0.30)^2}
=\SI{0.50}{\metre}.
\]

El potencial total es

\[
V
=
k
\left(
\frac{q_1}{r_1}
+
\frac{q_2}{r_2}
\right).
\]

Sustituyendo,

\[
V
=
(8.99\times10^9)
\left[
\frac{3.0\times10^{-6}}{0.30}
-
\frac{4.0\times10^{-6}}{0.50}
\right].
\]

Entonces,

\[
V
=
(8.99\times10^9)
(1.0\times10^{-5}-8.0\times10^{-6}).
\]

Por tanto,

\[
\boxed{
V\approx\SI{1.8e4}{\volt}
}.
\]

Observe que no fue necesario hallar direcciones ni componentes: el potencial es una
magnitud escalar.


\section{Puntos de potencial nulo}

Para varias cargas puede existir un punto donde

\[
V=0.
\]

Para dos cargas,

\[
k\frac{q_1}{r_1}
+
k\frac{q_2}{r_2}
=0.
\]

Cancelando \(k\),

\begin{equation}
\boxed{
\frac{q_1}{r_1}
+
\frac{q_2}{r_2}
=0
}.
\end{equation}

Para que exista cancelación entre dos términos, sus signos deben ser opuestos.

Por ello, dos cargas de signos opuestos pueden producir puntos donde el potencial
total sea cero.


\begin{advertencia}

Un punto donde

\[
V=0
\]

no es necesariamente un punto donde

\[
E=0.
\]

El potencial es una suma escalar, mientras que el campo depende de la suma de sus
componentes.

Son condiciones físicamente diferentes.

\end{advertencia}


\subsection{Ejemplo conceptual: dipolo eléctrico}

Consideremos dos cargas

\[
+q
\qquad\text{y}\qquad
-q
\]

de igual magnitud.

En cualquier punto de la mediatriz del segmento que las une, las distancias a ambas
cargas son iguales:

\[
r_+=r_-.
\]

Entonces,

\[
V
=
k\frac{q}{r_+}
+
k\frac{-q}{r_-}
=0.
\]

Por tanto, toda la mediatriz es una región de potencial nulo.

Sin embargo, los campos eléctricos producidos por ambas cargas no se cancelan en
general.

Así,

\[
\boxed{
V=0
\not\Rightarrow
E=0
}.
\]


\section{Potencial eléctrico de una distribución continua de carga}

Cuando la carga está distribuida continuamente, dividimos la distribución en elementos
infinitesimales \(dq\).

Cada elemento produce

\[
dV=k\frac{dq}{r},
\]

donde \(r\) es la distancia entre \(dq\) y el punto donde se calcula el potencial.

El potencial total es

\begin{equation}
\boxed{
V=k\int\frac{dq}{r}
}.
\end{equation}

Esta integral es escalar.


\subsection{Distribución lineal}

Si

\[
dq=\lambda\,d\ell,
\]

entonces

\begin{equation}
\boxed{
V=k\int\frac{\lambda\,d\ell}{r}
}.
\end{equation}


\subsection{Distribución superficial}

Si

\[
dq=\sigma\,dA,
\]

entonces

\begin{equation}
\boxed{
V=k\int\frac{\sigma\,dA}{r}
}.
\end{equation}


\subsection{Distribución volumétrica}

Si

\[
dq=\rho\,dV',
\]

entonces

\begin{equation}
\boxed{
V=k\int\frac{\rho\,dV'}{r}
}.
\end{equation}

Se utiliza \(dV'\) para representar aquí un elemento de volumen y evitar confundirlo
con el símbolo \(V\) del potencial eléctrico.


\section{Ventaja del potencial en distribuciones continuas}

Para calcular el campo eléctrico de una distribución continua suele ser necesario
analizar direcciones y componentes de cada contribución.

Para calcular el potencial sólo se suman cantidades escalares:

\[
dV=k\frac{dq}{r}.
\]

Por ello, en muchas geometrías es más sencillo obtener primero \(V\) y después
determinar el campo a partir de la variación espacial del potencial.


\section{Referencia del potencial}

Al igual que la energía potencial, el potencial eléctrico requiere una referencia.

Podemos sumar una constante \(C\) a todos los potenciales,

\[
V'=V+C,
\]

sin modificar ninguna diferencia de potencial:

\[
V'_B-V'_A
=
V_B-V_A.
\]

Por tanto, las predicciones físicas dependen de

\[
\Delta V,
\]

no del valor elegido arbitrariamente como cero.


\subsection{Referencia en el infinito}

Para distribuciones de carga localizadas y finitas suele elegirse

\[
\boxed{
V(\infty)=0
}.
\]

Esta elección conduce, por ejemplo, a

\[
V=k\frac{Q}{r}
\]

para una carga puntual.


\subsection{Distribuciones infinitas}

Para ciertas distribuciones idealizadas infinitas no es posible utilizar
\(V(\infty)=0\) de manera útil.

Esto ocurre, por ejemplo, con una línea infinita de carga o una lámina infinita.

En estos casos se escoge un punto de referencia finito y se trabaja directamente con
la diferencia de potencial

\[
V_B-V_A.
\]

\begin{idea}

El valor cero del potencial es una elección de referencia.

La diferencia de potencial entre dos puntos es la cantidad físicamente relevante.

\end{idea}


\section{Potencial y energía potencial}

La relación

\[
U=qV
\]

permite conectar el campo creado por unas cargas fuente con la energía de otra carga
que se coloca en él.

Si una carga \(q\) se mueve entre \(A\) y \(B\),

\[
\Delta U=q(V_B-V_A).
\]

Por tanto,

\begin{equation}
\boxed{
\Delta U=q\Delta V
}.
\end{equation}

Esta relación debe utilizar el signo real de \(q\).


\subsection{Carga positiva}

Si

\[
q>0,
\]

entonces \(U\) y \(V\) cambian en el mismo sentido:

\[
\Delta V>0
\Longrightarrow
\Delta U>0.
\]


\subsection{Carga negativa}

Si

\[
q<0,
\]

entonces \(U\) y \(V\) cambian en sentidos opuestos:

\[
\Delta V>0
\Longrightarrow
\Delta U<0.
\]

Esta diferencia explica por qué una carga negativa tiende a moverse espontáneamente
hacia potenciales mayores.


\section{Conservación de la energía usando el potencial}

Si sólo intervienen fuerzas electrostáticas,

\[
K_i+U_i=K_f+U_f.
\]

Como

\[
U=qV,
\]

podemos escribir

\begin{equation}
\boxed{
K_i+qV_i=K_f+qV_f
}.
\end{equation}

Por tanto,

\begin{equation}
\boxed{
K_f-K_i
=
-q(V_f-V_i)
}.
\end{equation}

Es decir,

\begin{equation}
\boxed{
\Delta K=-q\Delta V
}.
\end{equation}


\subsection{Partícula liberada desde el reposo}

Si

\[
K_i=0,
\]

entonces

\[
K_f=q(V_i-V_f).
\]

Como

\[
K_f=\frac12mv_f^2,
\]

se obtiene

\begin{equation}
\boxed{
\frac12mv_f^2
=
q(V_i-V_f)
}.
\end{equation}

Esta expresión sólo puede utilizarse si el lado derecho produce una energía cinética
no negativa.


\section{Superficies equipotenciales}

Una \textbf{superficie equipotencial} es una superficie cuyos puntos tienen el mismo
potencial eléctrico.

Si dos puntos \(A\) y \(B\) pertenecen a una misma superficie equipotencial,

\[
V_A=V_B.
\]

Por tanto,

\[
\boxed{\Delta V=0}.
\]

Como

\[
\Delta U=q\Delta V,
\]

también

\[
\boxed{\Delta U=0}.
\]

Además,

\[
W_e=-\Delta U,
\]

por lo que

\begin{equation}
\boxed{
W_e=0
}
\end{equation}

para cualquier desplazamiento realizado completamente sobre una superficie
equipotencial.


\subsection{Equipotenciales de una carga puntual}

Para una carga puntual,

\[
V=k\frac{Q}{r}.
\]

Todos los puntos que se encuentran a la misma distancia \(r\) tienen el mismo
potencial.

Por tanto, las superficies equipotenciales son esferas concéntricas centradas en la
carga.

En una representación bidimensional aparecen como circunferencias concéntricas.

\begin{center}
\begin{tikzpicture}[>=Latex,scale=0.9]

\filldraw (0,0) circle (2.5pt)
node[below right=3pt] {$+Q$};

\draw (0,0) circle (1);
\draw (0,0) circle (2);
\draw (0,0) circle (3);

\foreach \ang in {0,45,...,315}
{
    \draw[->,thick]
    ({0.35*cos(\ang)},{0.35*sin(\ang)})
    --
    ({3.5*cos(\ang)},{3.5*sin(\ang)});
}

\node at (0.7,0.7) {$V_1$};
\node at (1.4,1.4) {$V_2$};
\node at (2.15,2.15) {$V_3$};

\end{tikzpicture}
\end{center}

Para una carga positiva,

\[
V_1>V_2>V_3
\]

si las superficies están ordenadas desde la más cercana hasta la más lejana.


\section{Campo eléctrico y superficies equipotenciales}

Las líneas de campo eléctrico son siempre perpendiculares a las superficies
equipotenciales.

Para comprenderlo, consideremos un pequeño desplazamiento \(d\ell\) realizado sobre
una superficie equipotencial.

Como

\[
dV=0,
\]

se cumple

\[
E_x\,dx+E_y\,dy+E_z\,dz=0.
\]

Esto significa que el campo no posee componente en la dirección tangente a la
superficie.

Por tanto,

\begin{equation}
\boxed{
\vec E\perp\text{superficie equipotencial}
}.
\end{equation}

El campo apunta además desde valores mayores de \(V\) hacia valores menores de \(V\).


\begin{idea}

Las dos reglas fundamentales son

\[
\boxed{
\text{las líneas de campo son perpendiculares a las equipotenciales}
}
\]

y

\[
\boxed{
\vec E\text{ apunta hacia potencial decreciente}
}.
\]

\end{idea}


\section{Separación entre superficies equipotenciales}

La separación entre superficies equipotenciales proporciona información sobre la
magnitud del campo.

Si el potencial cambia mucho en una distancia pequeña, el campo es intenso.

Si cambia lentamente con la posición, el campo es menor.

Por tanto,

\[
\boxed{
\text{equipotenciales muy próximas}
\Longrightarrow
\text{campo intenso}
}.
\]

En cambio,

\[
\boxed{
\text{equipotenciales muy separadas}
\Longrightarrow
\text{campo menos intenso}
}.
\]


\section{Obtención del campo eléctrico a partir del potencial}

En una dimensión,

\[
V_B-V_A=-\int_A^B E_x\,dx.
\]

Para un desplazamiento infinitesimal,

\[
dV=-E_x\,dx.
\]

Por tanto,

\begin{equation}
\boxed{
E_x=-\frac{dV}{dx}
}.
\end{equation}

El campo es el negativo de la rapidez con que cambia el potencial con la posición.


\subsection{Tres dimensiones}

Si el potencial depende de las coordenadas

\[
V=V(x,y,z),
\]

las componentes del campo son

\begin{equation}
\boxed{
E_x=-\frac{\partial V}{\partial x}
}
\end{equation}

\[
\boxed{
E_y=-\frac{\partial V}{\partial y}
}
\]

\[
\boxed{
E_z=-\frac{\partial V}{\partial z}
}.
\]

Estas expresiones son especialmente importantes porque permiten obtener directamente
las componentes rectangulares del campo a partir de una función escalar \(V(x,y,z)\).


\subsection{Interpretación gráfica}

En una gráfica de \(V\) frente a \(x\),

\[
E_x=-\frac{dV}{dx}.
\]

Por tanto:

\[
\frac{dV}{dx}>0
\Longrightarrow
E_x<0,
\]

\[
\frac{dV}{dx}<0
\Longrightarrow
E_x>0,
\]

y

\[
\frac{dV}{dx}=0
\Longrightarrow
E_x=0.
\]

Una pendiente muy pronunciada corresponde a un campo de gran magnitud.


\section{Ejemplo: obtener el campo a partir del potencial}

\begin{ejemplo}

En cierta región del espacio el potencial depende de \(x\) según

\[
V(x)=50-20x+4x^2,
\]

donde \(V\) se expresa en voltios y \(x\) en metros.

Determine \(E_x\) en

\[
x=\SI{2.0}{\metre}.
\]

\end{ejemplo}

Utilizamos

\[
E_x=-\frac{dV}{dx}.
\]

Derivando,

\[
\frac{dV}{dx}=-20+8x.
\]

Por tanto,

\[
E_x=20-8x.
\]

En

\[
x=\SI{2.0}{\metre},
\]

obtenemos

\[
E_x=20-8(2.0).
\]

Así,

\[
\boxed{
E_x=+\SI{4.0}{\volt\per\metre}
}.
\]

Como

\[
\SI{1}{\volt\per\metre}
=
\SI{1}{\newton\per\coulomb},
\]

también puede escribirse

\[
\boxed{
E_x=+\SI{4.0}{\newton\per\coulomb}
}.
\]

El signo positivo indica que el campo apunta hacia \(+x\).


\section{Unidad alternativa del campo eléctrico}

A partir de

\[
E_x=-\frac{dV}{dx},
\]

la unidad del campo puede expresarse como

\[
\boxed{
[E]=\si{\volt\per\metre}
}.
\]

Esta unidad es equivalente a

\[
\boxed{
\SI{1}{\volt\per\metre}
=
\SI{1}{\newton\per\coulomb}
}.
\]


\section{Potencial de una línea infinita}

Para una línea infinita uniformemente cargada,

\[
E=\frac{|\lambda|}{2\pi\varepsilon_0r}.
\]

No puede elegirse de manera útil

\[
V(\infty)=0,
\]

porque la integral hasta el infinito no converge.

En cambio, podemos calcular la diferencia de potencial entre dos distancias
\(r_A\) y \(r_B\).

Para \(\lambda>0\),

\[
V_B-V_A
=
-\int_{r_A}^{r_B}
\frac{\lambda}{2\pi\varepsilon_0r}\,dr.
\]

Por tanto,

\begin{equation}
\boxed{
V_B-V_A
=
-\frac{\lambda}{2\pi\varepsilon_0}
\ln\left(\frac{r_B}{r_A}\right)
}.
\end{equation}

Esta expresión ilustra por qué, en ciertas distribuciones infinitas, debe trabajarse
con diferencias de potencial en lugar de asignar cero en el infinito.


\section{Potencial en un campo uniforme}

Si el campo uniforme está dirigido sobre el eje \(x\),

\[
E_x=\text{constante}.
\]

Entonces,

\[
dV=-E_x\,dx.
\]

Integrando,

\[
V(x)=-E_xx+C,
\]

donde \(C\) depende de la referencia elegida.

Así,

\begin{equation}
\boxed{
V(x)=V_0-E_x(x-x_0)
}.
\end{equation}

Por tanto, en un campo uniforme el potencial cambia linealmente con la posición.


\section{Campo nulo y potencial constante}

Si en toda una región

\[
E_x=E_y=E_z=0,
\]

entonces

\[
\frac{\partial V}{\partial x}
=
\frac{\partial V}{\partial y}
=
\frac{\partial V}{\partial z}
=0.
\]

Por tanto,

\[
\boxed{
E=0
\Longrightarrow
V=\text{constante en la región}
}.
\]

La constante no tiene que ser cero.


\begin{advertencia}

Las afirmaciones

\[
V=0
\]

y

\[
E=0
\]

no son equivalentes.

Puede existir un punto con

\[
V=0
\]

y

\[
E\neq0.
\]

También puede existir una región donde

\[
E=0
\]

pero

\[
V\neq0
\]

y sea simplemente constante.

\end{advertencia}


\section{Potencial y trayectoria}

Como el campo electrostático es conservativo,

\[
V_B-V_A
=
-\int_A^B
\left(
E_x\,dx+E_y\,dy+E_z\,dz
\right)
\]

es independiente de la trayectoria seguida entre \(A\) y \(B\).

Por tanto, si se conocen los puntos inicial y final, cualquier trayectoria conveniente
puede utilizarse para evaluar la diferencia de potencial.

En una trayectoria cerrada,

\[
V_f=V_i,
\]

de modo que

\begin{equation}
\boxed{
\oint
\left(
E_x\,dx+E_y\,dy+E_z\,dz
\right)=0
}.
\end{equation}


\section{Errores y confusiones frecuentes}

\begin{errorfrecuente}

\textbf{Confundir potencial eléctrico con energía potencial eléctrica.}

El potencial se mide en voltios:

\[
[V]=\si{\volt}.
\]

La energía potencial se mide en joules:

\[
[U]=\si{\joule}.
\]

Ambas magnitudes se relacionan mediante

\[
U=qV.
\]

\end{errorfrecuente}


\begin{errorfrecuente}

\textbf{Pensar que el potencial es un vector.}

El potencial eléctrico es escalar.

Por ello, para varias cargas se suman algebraicamente los potenciales:

\[
V=\sum_iV_i.
\]

\end{errorfrecuente}


\begin{errorfrecuente}

\textbf{Utilizar valor absoluto en el potencial de una carga puntual.}

La expresión correcta es

\[
V=k\frac{Q}{r}.
\]

Debe conservarse el signo de \(Q\).

\end{errorfrecuente}


\begin{errorfrecuente}

\textbf{Confundir el signo del potencial con la dirección del campo.}

El potencial no posee dirección.

La dirección del campo se obtiene de cómo cambia \(V\) con la posición.

\end{errorfrecuente}


\begin{errorfrecuente}

\textbf{Pensar que una carga positiva siempre se mueve hacia potencial mayor.}

Si se mueve espontáneamente bajo la acción del campo, una carga positiva tiende hacia
potenciales menores.

\end{errorfrecuente}


\begin{errorfrecuente}

\textbf{Aplicar la misma regla a una carga negativa.}

Como

\[
U=qV,
\]

una carga negativa disminuye su energía potencial al desplazarse hacia potenciales
mayores.

\end{errorfrecuente}


\begin{errorfrecuente}

\textbf{Pensar que \(V=0\) implica \(E=0\).}

No existe tal implicación.

Por ejemplo, en la mediatriz de un dipolo puede cumplirse

\[
V=0
\]

mientras el campo es diferente de cero.

\end{errorfrecuente}


\begin{errorfrecuente}

\textbf{Pensar que \(E=0\) implica necesariamente \(V=0\).}

Si el campo es cero en una región, el potencial es constante, pero dicha constante
puede ser distinta de cero.

\end{errorfrecuente}


\begin{errorfrecuente}

\textbf{Medir el potencial absoluto sin establecer una referencia.}

El potencial se define respecto de un nivel de referencia.

Para distribuciones localizadas suele elegirse

\[
V(\infty)=0.
\]

\end{errorfrecuente}


\begin{errorfrecuente}

\textbf{Usar \(V(\infty)=0\) indiscriminadamente para distribuciones infinitas.}

Para una línea infinita o una lámina infinita, esta elección no conduce a un potencial
finito útil.

Debe trabajarse con diferencias de potencial respecto de una referencia finita.

\end{errorfrecuente}


\begin{errorfrecuente}

\textbf{Olvidar el signo negativo en la relación entre campo y potencial.}

En una dimensión,

\[
E_x=-\frac{dV}{dx}.
\]

El campo apunta hacia donde el potencial disminuye.

\end{errorfrecuente}


\begin{errorfrecuente}

\textbf{Suponer que moverse sobre una superficie equipotencial requiere trabajo
eléctrico.}

Sobre una equipotencial,

\[
\Delta V=0,
\]

por lo que

\[
W_e=-q\Delta V=0.
\]

\end{errorfrecuente}


\section{Preguntas para comprobar comprensión}

\begin{enumerate}
    \item ¿Cómo se define el potencial eléctrico?

    \item ¿Cuál es la diferencia conceptual entre \(U\) y \(V\)?

    \item ¿Por qué el potencial eléctrico no depende de la carga de prueba?

    \item ¿Es el potencial una magnitud escalar o vectorial?

    \item ¿Cuál es la unidad SI del potencial?

    \item ¿Qué representa físicamente \(\SI{1}{\volt}\)?

    \item ¿Cómo se define la diferencia de potencial?

    \item ¿Qué relación existe entre \(\Delta U\) y \(\Delta V\)?

    \item ¿Qué relación existe entre el trabajo de la fuerza eléctrica y
    \(\Delta V\)?

    \item ¿Hacia potenciales mayores o menores apunta el campo eléctrico?

    \item ¿Hacia qué potencial tiende a desplazarse espontáneamente una carga positiva?

    \item ¿Y una carga negativa?

    \item ¿Cuál es el potencial producido por una carga puntual?

    \item ¿Por qué no debe utilizarse valor absoluto en
    \[
    V=k\frac{Q}{r}?
    \]

    \item ¿Cómo se calcula el potencial producido por varias cargas puntuales?

    \item ¿Por qué el cálculo del potencial suele ser más sencillo que el del campo?

    \item ¿Puede existir \(V=0\) con \(E\neq0\)?

    \item ¿Puede existir \(E=0\) con \(V\neq0\)?

    \item ¿Qué es una superficie equipotencial?

    \item ¿Cuánto trabajo realiza el campo al desplazar una carga sobre una
    equipotencial?

    \item ¿Cómo se orientan las líneas de campo respecto de las superficies
    equipotenciales?

    \item ¿Qué indica una gran proximidad entre superficies equipotenciales?

    \item ¿Cómo se obtiene \(E_x\) a partir de una función \(V(x)\)?

    \item ¿Cómo se obtienen \(E_x\), \(E_y\) y \(E_z\) a partir de \(V(x,y,z)\)?

    \item ¿Por qué el potencial en un campo uniforme varía linealmente con la posición?

    \item ¿Por qué no siempre puede tomarse \(V(\infty)=0\)?

    \item ¿Qué significa físicamente elegir un nivel cero diferente para el potencial?

    \item ¿Depende la diferencia de potencial de la trayectoria seguida entre dos
    puntos?
\end{enumerate}


\begin{resumen}

\begin{itemize}
    \item El potencial eléctrico es energía potencial por unidad de carga:
    \[
    \boxed{V=\frac{U}{q}}.
    \]

    \item La energía potencial de una carga en un punto de potencial \(V\) es
    \[
    \boxed{U=qV}.
    \]

    \item La unidad SI es
    \[
    \boxed{\SI{1}{\volt}=\SI{1}{\joule\per\coulomb}}.
    \]

    \item La diferencia de potencial es
    \[
    \boxed{\Delta V=V_B-V_A}.
    \]

    \item Se relaciona con la energía potencial mediante
    \[
    \boxed{\Delta U=q\Delta V}.
    \]

    \item El trabajo eléctrico es
    \[
    \boxed{W_e=-q\Delta V}.
    \]

    \item La diferencia de potencial puede obtenerse del campo:
    \[
    \boxed{
    V_B-V_A
    =
    -\int_A^B
    (E_x\,dx+E_y\,dy+E_z\,dz)
    }.
    \]

    \item En un campo uniforme,
    \[
    \boxed{
    \Delta V
    =
    -(E_x\Delta x+E_y\Delta y+E_z\Delta z)
    }.
    \]

    \item Para una carga puntual, tomando \(V(\infty)=0\),
    \[
    \boxed{V=k\frac{Q}{r}}.
    \]

    \item Para varias cargas puntuales,
    \[
    \boxed{
    V=k\sum_i\frac{q_i}{r_i}
    }.
    \]

    \item Para una distribución continua,
    \[
    \boxed{
    V=k\int\frac{dq}{r}
    }.
    \]

    \item Una superficie equipotencial cumple
    \[
    \boxed{V=\text{constante}}.
    \]

    \item Sobre una superficie equipotencial,
    \[
    \boxed{\Delta V=0,\qquad W_e=0}.
    \]

    \item El campo eléctrico es perpendicular a las superficies equipotenciales y apunta
    hacia potenciales decrecientes.

    \item En una dimensión,
    \[
    \boxed{
    E_x=-\frac{dV}{dx}
    }.
    \]

    \item En coordenadas rectangulares,
    \[
    \boxed{
    E_x=-\frac{\partial V}{\partial x},
    \qquad
    E_y=-\frac{\partial V}{\partial y},
    \qquad
    E_z=-\frac{\partial V}{\partial z}
    }.
    \]

    \item Las unidades
    \[
    \boxed{\si{\volt\per\metre}}
    \]
    y
    \[
    \boxed{\si{\newton\per\coulomb}}
    \]
    son equivalentes para el campo eléctrico.

    \item Si sólo actúan fuerzas electrostáticas,
    \[
    \boxed{
    K_i+qV_i=K_f+qV_f
    }.
    \]

    \item Para una distribución localizada suele utilizarse
    \[
    \boxed{V(\infty)=0},
    \]
    pero para ciertas distribuciones infinitas debe elegirse una referencia finita.

    \item Un potencial nulo no implica un campo nulo:
    \[
    \boxed{V=0\not\Rightarrow E=0}.
    \]
\end{itemize}

\end{resumen}